Die Sprachliche Endkorrektur ist der aufwändigste Schritt bei der
Fertigstellung einer Arbeit und wird sich vermutlich auch mehrere Durchgänge
benötigen. Nichtsdestotrotz sollte dieser Schritt keinesfalls hlabherzig
angegangen werden. Eine Häufung von Tippfehlern wirft beim Gutachter zurecht
die Frage auf, ob der Autor auch in inhaltlichen Belangen die notwendige
Gründlichkeit hat walten lassen. Es gibt mehrere Möglichkeiten die korrektur
auf Rechtschreibung, Grammatik und Ausdruck zu bewerkstelligen. In der Regel
führt eine Kombination aller Varianten zum Ziel.


\paragraph{Automatische Rechtschreibprüfung}

Die automatische Rechtschreibprüfung ist die einfachste Methode der
Rechtschreibkorrektur. Mit ihr lassen sich am besten Tippfehler erkennen.
Viele moderne \LaTeX-Editoren bieten entsprechende Funktionen an. Ein normaler
Texteditor könnte u.\,U. von \LaTeX-Kommandos verwirrt werden. Einige Editoren
erkennen nicht nicht Tippfehler, sondern auch grammatikalisch fragwürdige
Konstruktionen und/oder heben Wiederholungen und die Verwendungen überflüssiger
Füllwörter vor.

Sollten Sie Ihre Arbeit mit der Kommandozeile unter UNIX-artigen Systemen
verfassen, bietet die Software Aspell\footnote{\url{http://aspell.net/}} die
Möglichkeit \TeX-Dateien zu überprüfen. Der Kommandozeilenaufruf dazu lautet:

\begin{lstlisting}[escapechar=\!]
$ aspell -t -c !\lara{Dateiname}!
\end{lstlisting}

Schlechten Ausdruck oder Satzbau erkennen solche Werkzeuge allerdings nicht.


\paragraph{Manuelle Korrektur eines Probedrucks}

Fertigen Sie frühzeitig einen Probedruck von Kapiteln an, die bereits (bis auf
Kleinigkeiten) abgeschlossen sind. Mitunter versperren \LaTeX-Kommandos im
Manuskript einen direkten Blick auf den fertigen Text, sodass der Blick auf das
fertige Ergebnis hilft, die Texte nun auch auf Satzbau und Grammatik
korrigieren zu können.

Tipp- und Rechtschreibfehler erkennen Sie leichter, wenn Sie jede Seite
rückwärts lesen. Sie vermeiden es so, Wörter beim Lesen zu überspringen oder
aus dem Lesefluss heraus anders zu lesen, als sie auf Papier stehen.

Streichen Sie erkannte Fehler sofort an. Streichen Sie die Fehlermarkierungen
erst dann in einer anderen Farbe aus, wenn der Fehler im \TeX-File behoben
wurde. Schmeißen Sie eine Seite des Probedrucks erst weg, wenn alle Fehler
ausgestrichen sind.


\paragraph{Korrektur durch unbeteiligte Bekannte}

Fragen Sie in Ihrem Bekanntenkreis, ob jemand bereit wäre, die Arbeit korrektur
zu lesen. Es gibt erfahrungsgemäß immer jemanden, der jemand kennt, der sich
mit den noch so kleinsten Feinheiten der Deutschen Sprache auskennt. Es kann
sogar von Hilfe sein, wenn der Korrekturleser nicht vom Fach ist. Er oder sie
wird sich dann sogar unweigerlich auf die Rechtschreibung und weniger auf den
Inhalt konzentrieren.

Um Anmerkungen schneller in das Manuskript einpflegen zu können, empfiehlt es
sich die Zeilen im Probedruck durchzunummierieren. Es ist viel leichter das
fünfte~Wort in Zeile~256 zu korrigieren, als Wort~fünf in der vierten~Zeile
von Absatz~3 im Abschnitt~2.3. \LaTeX\ fügt vollautomatisch Zeilennummern ein,
indem Sie in der Vorlage die in Listing~\ref{lst:final:lineno} angegebenen
Zeilen auskommentieren.

\lstinputlisting[language={[LaTeX]{TeX}},style=block,escapechar=\!,float={htb},caption={Einschalten der Zeilennummerierung},label={lst:final:lineno}]{code/final_lineno.tex}

Vergessen Sie nicht, die Kommentarzeichen vor Abgabe der Arbeit wieder
einzufügen.


\paragraph{Feedback des Betreuers}

Wenn es sich um eine Abschlussarbeit handelt, wird ein guter Betreuer eine
Vorabversion anfordert. Ja nach Arbeitsbelastung und Stellensituation können
hier im realen Leben Abweichungen zu beobachten sein. Mit etwas
Fingerspitzengefühl können Sie jedoch im Zweifel einen Betreuer dazu überreden
Ihre Arbeit zu lesen. Bedenken Sie aber, dass die Mühlen an einer Universität
langsam mahlen. Ihr Betreuer hat sicherlich parallel noch Anträge zu
unterschreiben, Lehrveranstaltungen zu halten und Konferenzen zu besuchen. Der
Betreuer ist aber nicht Ihr Sekretär. Tipp- und Ausdrucksfehler wird er Ihnen
sicher mitteilen, Fragen zu Fachtermini und zur inhaltlichen Schlüssigkeit
stehen hier klar im Vordergrund.
